% problem-000042 2 硫化氢沉淀与天平平衡
\section*{2. 硫化氢沉淀与天平平衡}
假定本题两杯中的液体体积始终等于1.00 L，且H_{2}S饱和溶液的浓度均等于 $0 . 1 0 0 \mathrm { m o l } \mathrm { L } _ { \mathrm { ~ } \circ } ^ { - 1 }$ 天平两盘各有⼀盛着 $1 . 0 \mathrm { { L H _ { 2 } O } }$ 的烧杯，天平平衡。

\subsection*{2-1}
向左杯中融⼊ $2 . 7 2 \ : \mathrm { g } \ : \mathrm { Z n C l _ { 2 } }$ ，试通过计算说明，为使天平再次平衡需向右杯中加⼊多少克 $\mathrm { N a _ { 2 } O _ { 2 } ? }$

\subsection*{2-2}
天平再次平衡后，向左杯中通⼊ $\mathrm { H } _ { 2 } \mathrm { S }$ ⽓体，充分反应且使 $\mathrm { H } _ { 2 } \mathrm { S }$ 达到饱和。接着向右杯中通⼊ $\mathrm { H } _ { 2 } \mathrm { S }$ ⽓体。试通过计算求出当天平第三次平衡时，右杯中溶液的 $\mathrm { p H } _ { \cdot }$ o

\subsection*{2-3}
天平第三次平衡后，向左杯中缓缓通⼊ $\mathrm { H C l }$ ⽓体。试通过计算求出当杯中的沉淀刚好溶解掉⼀半时，溶液的Cl−浓度。

\subsection*{2-4}
拟使天平第四次平衡，需向右杯中通⼊多少克HCl⽓体？试通过计算加以说明。
